\NormalFont\ShortTitle{Iacov}

{\MT Scrisoarea de la Iacov


\par }\ChapOne{1}{\PP \VerseOne{1}Iacov, rob al lui Dumnezeu și al Domnului Isus Hristos, către cele douăsprezece seminții care sunt în Dispersie: Salutări.


\par }{\PP \VS{2}Frații mei, socotiți că este o bucurie, când cădeți în diferite ispite,

\VS{3}știind că încercarea credinței voastre produce răbdare.

\VS{4}Lăsați ca răbdarea să-și facă lucrarea sa desăvârșită, ca să fiți desăvârșiți și compleți, fără să vă lipsească nimic.


\par }{\PP \VS{5}Iar dacă vreunul din voi are lipsă de înțelepciune, să ceară de la Dumnezeu, care dă tuturor cu generozitate și fără prihană, și i se va da.

\VS{6}Dar să ceară cu credință, fără să se îndoiască, căci cel ce se îndoiește este ca un val al mării, împins de vânt și zvârlit.

\VS{7}Căci acel om nu trebuie să creadă că va primi ceva de la Domnul.

\VS{8}El este un om cu minte dublă, instabil în toate căile sale.


\par }{\PP \VS{9}Fratele care se află într-o situație umilă să se laude cu poziția lui înaltă;

\VS{10}iar cel bogat, să se smerească, pentru că va trece ca floarea din iarbă.

\VS{11}Căci soarele răsare cu vânt arzător și usucă iarba; floarea din ea cade și frumusețea înfățișării ei piere. Așa va dispărea și bogatul în activitățile sale.


\par }{\PP \VS{12}Ferice de cel ce rabdă ispita, căci, după ce va fi fost învrednicit, va primi cununa vieții, pe care Domnul a făgăduit-o celor ce-L iubesc.


\par }{\PP \VS{13}Nimeni să nu zică, când este ispitit: “Sunt ispitit de Dumnezeu”, căci Dumnezeu nu poate fi ispitit de rău și El însuși nu ispitește pe nimeni.

\VS{14}Ci fiecare este ispitit atunci când este atras de poftele sale și ispitit.

\VS{15}Atunci pofta, după ce a conceput, naște păcatul. Păcatul, când a ajuns la maturitate, produce moartea.

\VS{16}Nu vă lăsați înșelați, iubiții mei frați.

\VS{17}Orice dar bun și orice dar desăvârșit vine de sus, coborând de la Tatăl luminilor, la care nu poate fi nici variație, nici umbră de întoarcere.

\VS{18}De bunăvoie ne-a născut prin cuvântul adevărului, ca să fim un fel de primii roade ai făpturilor sale.


\par }{\PP \VS{19}Așadar, iubiții mei frați, fiecare să fie iute la auz, încet la vorbă și încet la mânie;

\VS{20}căci mânia omului nu produce dreptatea lui Dumnezeu.

\VS{21}De aceea, lepădând orice murdărie și revărsare de răutate, primiți cu smerenie cuvântul implantat, care poate să vă mântuiască sufletele.


\par }{\PP \VS{22}Ci fiți făcători ai cuvântului, și nu numai ascultători, amăgindu-vă pe voi înșivă.

\VS{23}Căci, dacă cineva ascultă cuvântul și nu este făptuitor, se aseamănă cu un om care se uită în oglindă la chipul său natural;

\VS{24}căci se vede pe sine însuși, pleacă și îndată uită ce fel de om era.

\VS{25}Dar cel care privește în legea desăvârșită a libertății și continuă, nefiind un ascultător care uită, ci un făptuitor al lucrării, acesta va fi binecuvântat în ceea ce face.


\par }{\PP \VS{26}Dacă cineva dintre voi se crede credincios, dar nu-și înfrânează limba, ci își amăgește inima, atunci religia acestui om este fără valoare.

\VS{27}Religia curată și nepătată înaintea Dumnezeului și Tatălui nostru este aceasta: să vizitezi pe orfan și pe văduve în necazurile lor și să te păstrezi nepătat de lume.



\par }\Chap{2}{\PP \VerseOne{1}Frații mei, nu țineți cu părtinire credința Domnului nostru glorios Isus Hristos.

\VS{2}Căci, dacă în sinagoga voastră intră un om cu inel de aur, îmbrăcat în haine frumoase, și intră și un sărac îmbrăcat în haine murdare,

\VS{3}și voi îi acordați o atenție deosebită celui îmbrăcat în haine frumoase și îi spuneți: “Șezi aici, într-un loc bun”, iar celui sărac îi spuneți: “Stai acolo” sau “Șezi lângă scăunelul meu”,

\VS{4}nu cumva ați dat dovadă de părtinire între voi și nu cumva v-ați făcut judecători cu gânduri rele?

\VS{5}Ascultați, iubiții mei frați. Oare nu i-a ales Dumnezeu pe cei săraci în această lume pentru a fi bogați în credință și moștenitori ai Împărăției pe care a promis-o celor care îl iubesc?

\VS{6}Dar voi l-ați dezonorat pe cel sărac. Nu cumva cei bogați vă asupresc și nu vă târăsc personal în fața tribunalelor?

\VS{7}Nu blasfemiază ei numele onorabil cu care sunteți chemați?


\par }{\PP \VS{8}Dar dacă împlinești legea împărătească, după Scriptura: “Să iubești pe aproapele tău ca pe tine însuți”, bine faci.

\VS{9}Dar dacă dați dovadă de părtinire, săvârșiți un păcat, fiind condamnați de lege ca niște călcători de lege.

\VS{10}Căci oricine păzește întreaga Lege, dar se împiedică într-un singur punct, s-a făcut vinovat de toate.

\VS{11}Căci cel care a zis: “Să nu comiți adulter”, a zis și: “Să nu comiți crimă”. Or, dacă nu comiți adulter, dar săvârșești o crimă, ai devenit un călcător al legii.

\VS{12}Vorbește și acționează astfel ca niște oameni care vor fi judecați după legea libertății.

\VS{13}Căci judecata este fără milă pentru cel care nu a arătat milă. Mila triumfă asupra judecății.


\par }{\PP \VS{14}La ce bun, frații mei, dacă un om zice că are credință, dar nu are fapte? Oare credința îl poate mântui?

\VS{15}Și dacă un frate sau o soră este gol și nu are hrana zilnică,

\VS{16}iar unul dintre voi îi spune: “Du-te în pace! Încălziți-vă și săturați-vă!”, dar nu le-ați dat cele de care are nevoie trupul, la ce bun?

\VS{17}Tot așa și credința, dacă nu are fapte, este moartă în sine.

\VS{18}Da, un om va spune: “Tu ai credință, iar eu am fapte”. Arată-mi credința ta fără fapte, iar eu îți voi arăta credința mea prin faptele mele.


\par }{\PP \VS{19}Voi credeți că Dumnezeu este unul singur. Faceți bine. Demonii cred și ei — și tremură.

\VS{20}Dar vrei tu să știi, om deșert, că credința fără fapte este moartă?

\VS{21}Nu cumva Avraam, tatăl nostru, nu a fost justificat prin fapte, prin faptul că a adus pe Isaac, fiul său, pe altar?

\VS{22}Vedeți că credința a lucrat cu faptele lui și prin fapte s-a desăvârșit credința.

\VS{23}Astfel s-a împlinit Scriptura care zice: “Avraam a crezut pe Dumnezeu și i s-a socotit ca o dreptate” și a fost numit prietenul lui Dumnezeu.

\VS{24}Vedeți deci că prin fapte omul este îndreptățit, și nu numai prin credință.

\VS{25}În același fel, nu a fost oare și Rahab, prostituata, justificată prin fapte, atunci când a primit mesagerii și i-a trimis pe o altă cale?

\VS{26}Căci, după cum trupul separat de duh este mort, tot așa și credința separată de fapte este moartă.



\par }\Chap{3}{\PP \VerseOne{1}Fraților, să nu fiți mulți dintre voi învățători, știind că vom primi o judecată mai grea.

\VS{2}Căci toți ne poticnim în multe lucruri. Oricine nu se poticnește în cuvânt este o persoană desăvârșită, capabilă să țină în frâu și tot trupul.

\VS{3}Într-adevăr, noi punem mânere în gura cailor pentru ca ei să ne asculte și le îndrumăm tot corpul.

\VS{4}Iată că și corăbiile, deși sunt atât de mari și sunt conduse de vânturi năprasnice, sunt totuși ghidate de o cârmă foarte mică, oriunde dorește pilotul.

\VS{5}Tot așa și limba este un membru mic și se laudă cu lucruri mari. Vedeți cum un foc mic se poate extinde până la o pădure mare!

\VS{6}Și limba este un foc. Lumea nelegiuirii dintre mădularele noastre este limba, care spurcă tot trupul și dă foc cursului firii, și este incendiată de Gheenă.

\VS{7}Căci orice fel de animal, pasăre, târâtoare și creatură marină este îmblânzită și a fost îmblânzită de oameni;

\VS{8}dar nimeni nu poate îmblânzi limba. Ea este un rău neliniștit, plin de otravă mortală.

\VS{9}Cu ea îl binecuvântăm pe Dumnezeul și Tatăl nostru, iar cu ea îi blestemăm pe oamenii care sunt făcuți după chipul lui Dumnezeu.

\VS{10}Din aceeași gură iese binecuvântarea și blestemul. Frații mei, aceste lucruri nu ar trebui să fie așa.

\VS{11}Un izvor trimite oare din aceeași deschizătură apă dulce și apă amară?

\VS{12}Poate un smochin, frații mei, să dea măsline, sau o viță de vie smochine? Așadar, niciun izvor nu produce în același timp apă sărată și apă dulce.


\par }{\PP \VS{13}Cine este înțelept și priceput între voi? Să arate prin buna lui purtare că faptele lui sunt făcute cu blândețea înțelepciunii.

\VS{14}Dar dacă ai în inima ta gelozie amară și ambiție egoistă, nu te lăuda și nu minți împotriva adevărului.

\VS{15}Această înțelepciune nu este cea care coboară de sus, ci este pământească, senzuală și demonică.

\VS{16}Căci unde sunt gelozia și ambiția egoistă, acolo este confuzie și orice faptă rea.

\VS{17}Dar înțelepciunea care vine de sus este mai întâi curată, apoi pașnică, blândă, rezonabilă, plină de milă și de roade bune, fără părtinire și fără ipocrizie.

\VS{18}Or, rodul dreptății este semănat în pace de către cei care fac pace.



\par }\Chap{4}{\PP \VerseOne{1}De unde vin războaiele și luptele dintre voi? Nu vin ele din plăcerile voastre care se războiesc în mădularele voastre?

\VS{2}Voi poftiți, și nu aveți. Ucideți și poftiți, și nu puteți obține. Vă luptați și faceți război. Nu aveți, pentru că nu cereți.

\VS{3}Cereți, și nu primiți, pentru că cereți cu motive greșite, ca să cheltuiți pe plăcerile voastre.

\VS{4}Voi, adulterinilor și adulterinelor, nu știți că prietenia cu lumea este ostilitate față de Dumnezeu? De aceea, oricine vrea să fie prieten cu lumea se face dușman al lui Dumnezeu.

\VS{5}Sau credeți că Scriptura spune în zadar: “Duhul care trăiește în noi tânjește cu gelozie”?

\VS{6}Dar el dă mai mult har. De aceea se spune: “Dumnezeu se împotrivește celor mândri, dar dă har celor smeriți”.

\VS{7}Fiți deci supuși lui Dumnezeu. Împotriviți-vă diavolului și el va fugi de voi.

\VS{8}Apropiați-vă de Dumnezeu, și el se va apropia de voi. Curățați-vă mâinile, păcătoșilor. Curățați-vă inimile, voi, cei cu minte dublă.

\VS{9}Plângeți, jeliți și plângeți. Râsul vostru să se transforme în jale și bucuria voastră în tristețe.

\VS{10}Umiliți-vă în fața Domnului, și el vă va înălța.


\par }{\PP \VS{11}Fraților, nu vorbiți unul împotriva altuia. Cine vorbește împotriva unui frate și judecă pe fratele său, vorbește împotriva legii și judecă legea. Dar dacă judeci legea, nu ești un împlinitor al legii, ci un judecător.

\VS{12}Numai unul singur este făcătorul legii, care poate să mântuiască și să distrugă. Dar cine ești tu să judeci pe altul?


\par }{\PP \VS{13}Veniți acum, voi, care ziceți: “Astăzi sau mâine să mergem în cetatea aceasta, să petrecem acolo un an, să facem negoț și să câștigăm.”

\VS{14}Dar voi nu știți cum va fi viața voastră mâine. Căci ce este viața voastră? Căci voi sunteți un abur care apare pentru puțin timp și apoi dispare.

\VS{15}Căci ar trebui să spuneți: “Dacă Domnul vrea, vom trăi amândoi și vom face una sau alta.”

\VS{16}Dar acum vă lăudați cu mândria voastră. Orice astfel de laudă este rea.

\VS{17}Așadar, pentru cel care știe să facă binele și nu-l face, pentru el este păcat.



\par }\Chap{5}{\PP \VerseOne{1}Veniți acum, bogaților, plângeți și urlați din cauza nenorocirilor care se abat asupra voastră.

\VS{2}Bogățiile voastre sunt corupte și hainele voastre sunt mâncate de molii.

\VS{3}Aurul și argintul vostru sunt corodate, iar coroziunea lor va fi ca o mărturie împotriva voastră și vă va mânca carnea ca focul. Voi v-ați pus comoara în zilele din urmă.

\VS{4}Iată că salariile muncitorilor care au secerat câmpurile voastre, pe care le-ați reținut prin fraudă, strigă; și strigătele celor care au secerat au intrat în urechile Domnului oștirilor.

\VS{5}Ați trăit în lux pe pământ și v-ați făcut plăcere. V-ați hrănit inimile ca într-o zi de măcel.

\VS{6}Ați condamnat și ați ucis pe cel neprihănit. El nu vă rezistă.


\par }{\PP \VS{7}Fiți răbdători, fraților, până la venirea Domnului. Iată că agricultorul așteaptă rodul prețios al pământului, răbdând asupra lui, până când primește ploaia timpurie și târzie.

\VS{8}Aveți și voi răbdare. Stabiliți-vă inimile, căci venirea Domnului este aproape.


\par }{\PP \VS{9}Fraților, nu vă certați unii pe alții, ca să nu fiți judecați. Iată că judecătorul stă la ușă.

\VS{10}Luați, fraților, ca exemplu de suferință și de perseverență, pe profeții care au vorbit în numele Domnului.

\VS{11}Iată, noi îi numim fericiți pe cei care au răbdat. Ați auzit despre perseverența lui Iov și ați văzut cum Domnul a ajuns la rezultat și cum Domnul este plin de compasiune și de îndurare.


\par }{\PP \VS{12}Dar mai presus de toate, frații mei, să nu jurați, nici pe cer, nici pe pământ, nici pe vreun alt jurământ, ci să fie “da” și “nu”, “nu”, ca să nu cădeți în ipocrizie.


\par }{\PP \VS{13}Suferă cineva dintre voi? Să se roage. Este vreunul dintre ei vesel? Să cânte laude.

\VS{14}Este cineva dintre voi bolnav? Să cheme pe bătrânii adunării și să se roage pentru el, ungându-l cu untdelemn în Numele Domnului;

\VS{15}și rugăciunea credinței îl va vindeca pe cel bolnav și Domnul îl va învia. Dacă a săvârșit păcate, va fi iertat.

\VS{16}Mărturisiți-vă păcatele unul altuia și rugați-vă unul pentru altul, ca să fiți vindecați. Rugăciunea insistentă a unei persoane neprihănite este puternic eficientă.

\VS{17}Ilie era un om cu o fire asemănătoare cu a noastră și s-a rugat cu insistență ca să nu plouă, iar pe pământ nu a plouat timp de trei ani și șase luni.

\VS{18}El s-a rugat din nou și cerul a dat ploaie, iar pământul și-a produs roadele.


\par }{\PP \VS{19}Fraților, dacă vreunul dintre voi se abate de la adevăr și cineva îl întoarce înapoi,

\VS{20}să știe că cel ce întoarce pe un păcătos de la rătăcirea căii lui va mântui un suflet de la moarte și va acoperi o mulțime de păcate.


\par }